% !TEX root = preview.tex
% =============================================================================
% Benchmark + evaluation (paper body, compact)
% Detailed generator pseudocode: sec_benchmark_appendix.tex
% =============================================================================

\section{Benchmark Construction}
\label{sec:benchmark}

Classical MAPF suites do not capture permanent geometric occupancy or
multi-terminal pipe groups, and we are not aware of a public 3D multi-group
pipe-routing suite with controlled difficulty and a constructive feasibility
certificate. We therefore build instances synthetically.

Each instance is a triple $(\mathcal{E},\{G_i\}_{i=1}^{m},\mathcal{W})$:
environment $\mathcal{E}$, $m$ groups with face terminals and uniform radius
$r$, and a conflict-free \textbf{witness} $\mathcal{W}$ proving generator-level
feasibility. Solvers never see $\mathcal{W}$ as a hint; it is used for
validation and as a \emph{reference length}, not an optimality bound.
Growing $\mathcal{W}$ first, then painting obstacles in its complement, yields
known-feasible clutter without coupling the suite to any solver under test.

\paragraph{Pipeline.}
(1)~Grow tree-like witness corridors with \textsc{AxisRay} or
\textsc{PlaneSlice} (groups attach in round-robin).
(2)~Paint static obstacles outside thickened witnesses with
\textsc{PillarField} or \textsc{CorridorRack}, never invading clearance $r$ on
$\mathcal{W}$.
(3)~Emit YAML (grid, obstacles, terminals, witness).
The product of two corridor methods and two obstacle families gives
\textbf{four datasets}, each with levels
\texttt{easy}/\texttt{middle}/\texttt{hard}
(Table~\ref{tab:difficulty}; $N{=}100$ successes per cell, domain $32^3$,
$r{=}0.5$).

\begin{table}[t]
  \centering
  \caption{Difficulty schedule ($m$ groups $\times$ $n_t$ terminals per group).
  Shared: domain $32^3$, $r{=}0.5$.}
  \label{tab:difficulty}
  \begin{tabular}{@{}lll@{}}
    \toprule
    Level & \textsc{AxisRay} ($m \times n_t$) & \textsc{PlaneSlice} ($m \times n_t$) \\
    \midrule
    easy   & $3\times 3$ & $3\times 3$ \\
    middle & $6\times 6$ & $8\times 8$ \\
    hard   & $9\times 9$ & $12\times 12$ \\
    \bottomrule
  \end{tabular}
\end{table}

\paragraph{Corridor generators.}
Both produce multi-terminal trees, not a single $s$--$t$ path; inter-group
corridors are obstacles for later growth.
\textsc{AxisRay} (preferred) grows an axis-aligned polyline tree by LOS scans
(no subdomain grid, no A$^\star$), with path/branch gaps and forced axis
changes at turns.
\textsc{PlaneSlice} grows on axis-aligned planes with \emph{planar} A$^\star$
only (no 3D fallback), using an $N{\times}N{\times}N$ subdomain layout and a
plane-gap within each group.
Pseudocode and clearance rules appear in Appendix~\ref{app:generators}.

\paragraph{Obstacle generators.}
Witness cells are thickened into a keep-out set $C$; obstacles never invade
$C$.
\textsc{PillarField} places large AABBs until target density
$\rho^\star{=}V_{\mathrm{obs}}/V_0$ (we use $\rho^\star{=}0.4$), sampling seeds
by distance to blocked cells.
\textsc{CorridorRack} dilates witnesses by $k$ Chebyshev steps
($k{=}2$ in reported runs) and meshes the complement as AABBs (density derived).
Details in Appendix~\ref{app:generators}.

\section{Evaluation Protocol}
\label{sec:eval}

Given $(\mathcal{E},\{G_i\})$ without $\mathcal{W}$, a solver must return
conflict-free paths for all topology tasks or fail within budget.
Reported runs fix topology to pure MST and low-level search to A$^\star$
(26-connected, Manhattan), and compare CBS / PBS / PP as in
Sec.~\ref{sec:method}.

\paragraph{Primary metric.}
Success rate under wall-clock budget $T$ (easy $30\,$s, middle $60\,$s, hard
$120\,$s):
\[
  \mathrm{SuccessRate}
  = \frac{1}{N}\bigl|\{ i : \text{instance $i$ solved within $T$}\}\bigr|.
\]
We also report failure reasons (\texttt{timeout} / \texttt{no\_solution}) and
success-vs-time curves. Fixed-budget success matches batch design use better
than optimality claims we do not make.

\paragraph{Secondary metrics.}
Runtime on successes; on successes,
$\mathrm{cost\_ratio}=\mathrm{cost}(\mathrm{sol})/\mathrm{cost}(\mathcal{W})$
as a descriptive quality signal (\emph{not} an optimality gap---$\mathcal{W}$
is constructive). Cost comparisons across algorithms use the common success
set to limit selection bias.

\paragraph{Reproducibility.}
The suite, batch runners, and JSONL summaries are in the public repository
\url{https://github.com/bladesaber/MAPF_Pipeline}
(branch \texttt{branch\_quan}, commit \texttt{cf0005685626dcaedb4c96a64fdca554312dfba1}).
Reported runs use \texttt{pipe\_design/include/benchmark/gen\_robust.sh},
\texttt{run\_bench.sh}, and \texttt{pipe\_design/scripts/analyze\_runs.py}.
